\section{Conclusion}
\label{sec:conclusion}

We introduced a dual-constraint bisection framework for partisan-proportional
redistricting, extending minimum-edge-cut algorithms to $\mathtt{ncon}=2$ with
vertex weights $[\text{population},\, D_{\text{votes}}]$.
The central contribution is a closed-form formula~\eqref{eq:tpwgts-formula}
for the partition targets that guarantee proportional seat allocation:
\[
  \mathtt{tpwgts} = \left[\frac{k_D}{k},\; 1 - \frac{k_R}{2dk},\;
                          \frac{k_R}{k},\; \frac{k_R}{2dk}\right]
\]

The partisan Lorenz curve determines when this target is geographically feasible.
The partisan signal strength $\sigma$ measures how much the neutral geographic
distribution must be ``pulled'' toward the proportional target --- from
$\sigma \approx 0$ (free proportionality, geography already gives proportional
outcomes) to $\sigma > 0.2$ (expensive proportionality, geographic concentration
makes proportional maps nearly impossible).

\paragraph{The core finding: the proportionality paradox.}
Competitive states (46--55\,\% Democratic) have $\sigma < 0.05$ --- the
proportional tpwgts target is essentially the same as the neutral equal-distribution
target. Georgia needs to shift 0.1\,\% of D votes from neutral; Wisconsin needs
0.3\,\%. Yet both have 12--14\,pp gaps under geography-only algorithms.

By Theorem~\ref{thm:invariance} (Tradeoff invariance), varying the partisan
constraint from loose ($\eta = \infty$) to strict ($\eta = 1.0$) cannot close
more than $C \cdot \sigma$ of the gap for competitive states. For Georgia
($\sigma = 0.001$) and Wisconsin ($\sigma = 0.003$), this is less than 0.5\,pp
improvement on a 14\,pp gap.

\textbf{The proportionality problem in competitive states is a Lorenz feasibility
problem, not a target calibration problem.}
The algorithm cannot find the right partition because geography prevents it,
not because we are asking for the wrong target.

The states that most need proportional redistricting are the states where
proportionality is most expensive.
Extreme Democratic geographic concentration --- Milwaukee, Madison, Atlanta,
Boston --- produces large Rodden gaps \emph{and} makes proportional packing
geometrically difficult.
There is no geographic algorithm that can close the proportionality gap
in these states without explicitly incorporating partisan data.
This is a fact about geography.

\paragraph{What this means.}
For states with low $\sigma$ (cheap proportionality): geography-only algorithms
are nearly proportional. Small improvements to the geographic constraint
(B.9-style area balance or B.8-style isoperimetric normalisation) may close
the remaining gap without partisan data.

For states with high $\sigma$ (expensive proportionality): the proportionality
gap is a structural feature of the state's population distribution.
A court reviewing such a state's redistricting for partisan fairness should
understand that even a perfectly neutral algorithm produces a large gap ---
the gap is not evidence of manipulation, but of geography.

\paragraph{Future work.}
Three directions merit further investigation:
\begin{enumerate}
\item \textbf{Recursive proportionality}: applying the B.12 constraint at
  every level of the bisection tree to guarantee proportionality not just
  at the first split but throughout the entire district structure.

\item \textbf{Continuous tolerance curves}: empirical measurement of how the
  proportionality gap decreases as $\mathtt{ubvec}[1]$ is tightened, for each
  state in the 50-state sweep.

\item \textbf{Multi-member districts}: for expensive-proportionality states,
  characterising how many seats per district are needed before geographic
  constraints stop blocking proportionality --- connecting this analysis to
  the Huntington-Hill minimum from B.11.
\end{enumerate}
